\section{Reusable proof templates}

\subsection{Template: deterministic numerical operation}

\textbf{Data.}
PR terms \(f\) and \(d\), input validity formula \(V_X(x,d(q),o)\), output
validity formula \(V_Y(f(x),q,f_o(x,q,o))\).

\textbf{Local theorem.}
A \(Q\)-free native derivation of
\[
    V_X(x,d(q),o)\to V_Y(f(x),q,f_o(x,q,o)).
\]
Encode the implication as a clause or import its ordinary derivation.

\textbf{Graph use.}
The predecessor node exports input validity.  Specialize the local theorem at
the concrete terms.  List cut discharges the predecessor assumption.

\subsection{Template: guarded division}

\textbf{Data.}
Rational codes \(a,b\), a guard \(G(b,g)\) certifying
\(|b|\ge2^{-g}\), and requested precision \(q\).

\textbf{Demand.}
Choose \(r=q+c(g,M_a,M_b)\) so that input errors preserve
\(|b_r|\ge2^{-g-1}\) and the reciprocal error is below the budget.

\textbf{Verifier.}
Check denominator positivity, the rational lower bound, and the final cleared-
denominator error inequality.

\textbf{Output.}
A rational code \(a_r/b_r\) and an error certificate.  The totalized function
returns zero if the guard fails; the correctness clause has \(G\) as
antecedent.

\subsection{Template: bounded search with negative result}

\textbf{State.}
Clock, found flag, least witness, and Boolean prefix verifier.

\textbf{Invariant.}
If found, the witness is below the clock, satisfies the predicate, and all
smaller candidates fail.  If not found, every candidate below the clock fails.

\textbf{Output.}
A disjoint finite result: \path{found(w)} or \path{none}, each with its exact
contract.  Do not use a sentinel value without a status flag.

\subsection{Template: finite descent}

\textbf{Data.}
A fixed-dimensional bounded measure \(m(s)\), a step that either halts or
strictly lowers the measure, and a bound vector \(B\).

\textbf{Fuel.}
Use the mixed-radix state count \(\prod_i(B_i+1)\), or a sharper rank bound.

\textbf{Invariant.}
The current state is valid; the measure lies in the box; every nonhalt step
strictly lowers rank; once halted, the result is preserved.

\textbf{Conclusion.}
The trace halts within the fuel and the terminal certificate is valid.

\subsection{Template: cyclic precision contract}

\textbf{Data.}
A finite box \(L\), monotone demand transformer \(F\), and initial request
\(a\).

\textbf{Algorithm.}
Iterate \(x\mapsto x\join F(x)\) for
\(1+\sum_i(B_i-a_i)\) steps, recording a failure if \(F(x)\notin L\).

\textbf{Invariant.}
The sequence is monotone, remains in the box, and is below every prefixed point
above \(a\).

\textbf{Conclusion.}
Either a least closed demand is returned or a resource-coordinate overflow is
certified.

\subsection{Template: relational cloud}

\textbf{Data.}
Finite bounds \(B(q)\), validity \(G(q,w)\), and edge relation \(E(q,w,v)\).

\textbf{Algorithm.}
Build viability tables backward; select a least finite path only if a concrete
path is needed.

\textbf{Observer.}
Supply exact invariance, a diameter bound, or a motion-tail bound.  Export the
observer, not the implementation path.

\textbf{Exact labels.}
Add a finite label family and a positive separation certificate, then use
Theorem~\ref{thm:gap}.

\subsection{Template: certificate-only theorem}

\textbf{Producer interface.}
A code \(c\) and a bounded checker \(\mathsf{ok}(x,c)=1\).

\textbf{Soundness theorem.}
A local object derivation of
\[
       \mathsf{ok}(x,c)=1\to\Phi(x,\mathsf{out}(c)).
\]

\textbf{Completeness options.}
Either prove a PR producer \(p(x)\) with \(\mathsf{ok}(x,p(x))=1\), accept
\(c\) as explicit input/hypothesis, or state a bounded relation of all accepted
certificates.  Do not infer completeness from successful tests.

\subsection{Template: finite local-to-global argument}

\textbf{Data.}
A coded finite cover, local observations, overlap graph, and compatibility
clauses.

\textbf{Control.}
Use a spanning tree or finite nerve traversal.  At each edge, transport the
local invariant.  Cycles are checked by bounded commutativity conditions.

\textbf{Global output.}
A glued observation or stable quotient class.  A unique global representative
requires a separate uniqueness theorem; otherwise retain the equivalence class.

\subsection{Template: direct algebra plus control}

\textbf{Algebra phase.}
Use \path{ipra_ring} or a direct \(\Prf\;\PRAAx\) lemma to prove an open
identity.

\textbf{Import.}
Apply Theorem~\ref{thm:import} and store the resulting native formula as a node.

\textbf{Control phase.}
Use bounded cases, search, or a trace in \(\KPrf\), specializing the algebra
node as needed.

\textbf{Export.}
Apply Theorem~\ref{thm:export}, then use round-trip equivalence to present the
ordinary formula in its original syntax.
